\CatchFileBetweenTags{\AlphaInvVal}{constants.tex}{AlphaInvVal}
\CatchFileBetweenTags{\AlphaSVal}{constants.tex}{AlphasVal}
\CatchFileBetweenTags{\WeakAngleVal}{constants.tex}{WeakAngleVal}
\CatchFileBetweenTags{\HiggsVEVVal}{constants.tex}{HiggsVEVVal}
\CatchFileBetweenTags{\FermiConstVal}{constants.tex}{FermiConstVal}
\CatchFileBetweenTags{\HiggsMassVal}{constants.tex}{HiggsMassVal}
\CatchFileBetweenTags{\ElectronYukawaVal}{constants.tex}{ElectronYukawaVal}
\CatchFileBetweenTags{\JarlskogEq}{constants.tex}{JarlskogEq}
\CatchFileBetweenTags{\JarlskogVal}{constants.tex}{JarlskogVal}
\CatchFileBetweenTags{\WBosonMassVal}{constants.tex}{WBosonMassVal}
\CatchFileBetweenTags{\VacuumEnergyScaleVal}{constants.tex}{VacuumEnergyScaleVal}
\CatchFileBetweenTags{\combinedtotalchiVal}{constants.tex}{combinedtotalchiVal}
\CatchFileBetweenTags{\combinedreducedchiVal}{constants.tex}{combinedreducedchiVal}

\begin{abstract}
Persistence Mechanics constrains how a system can service fixed demands within finite resources, preserve distinctions under coarse descriptions, and transfer information between nested subsystems. We formalize capacity, reversible-encoding, reflected-load and nesting-depth conditions in Lean independently of the vacuum example. Applying these rules to the proposed $E_8$ substrate gives candidate allocation laws for gauge fractions, scalar regulation and bulk attenuation. Lean verifies conditional selection results, exact arithmetic and numerical enclosures; the substrate-specific workloads, response laws and physical matching remain explicit premises. Joint tests expose obligations not visible in separate numerical comparisons. The proposed VEV fixes the electron Yukawa through $m_e=y_e v/\sqrt{2}$; an independent dressing $y_e=PM(1+c\alpha)$ would need $c\in(1.28777,1.28779)$, which excludes the earlier $\sigma/D$ proposal. The mass ratio $m_W^2=(1-s)m_Z^2$ holds at every interface scale, but the absolute W and Z masses need one interface factor $k\in(1.0362,1.0365)$ on $\alpha$; the same factor cannot serve the strong readout, and with the measured Z mass it implies $\Delta r\approx0.0351$ and $m_W\approx80.385$~GeV, a falsifiable value distinct from standard radiative corrections. For the scalar interface, maximal gain cannot pay for its own reflection; the largest persistent gain lies within $c\cdot10^{-8}/23$ of $128/989$ for reflection charge rate $c$. Running the strong coupling with Paper~V's own flavour thresholds replaces a three-flavour, capacity-corrected estimate. These results provide quantitative tests for completing the substrate derivation. A unique physical determination of all constants, spectra and scale-dependent responses remains open.
\end{abstract}